%% kapa: exam=1 round=26 year=2015 period=1 form=A subject=civil_law subject_ko=민법 seq=24 number=24
%% source: Q-Net 감정평가사 기출문제 (cst00309) artlSeq=5138001 file=extracted/1008/5138001/제26회_감정평가사 1차_기출/제26회_감정평가사 1차_기출/1교시_감평_A형.pdf
%% conversion: kordoc 4.24.51 -> tools/gen1.py (auto); review=proofread-sonnet; flags=none
\begin{question}{24}
1990년 乙이 甲으로부터 토지를 매수하여 등기는 이전받지 아니한 채 인도받고 점유ㆍ사용하다가, 2003년 이를 丙이 乙로부터 매수하여 이전등기 없이 인도받고 점유ㆍ사용하고 있다. 옳은 설명을 모두 고른 것은? (다툼이 있으면 판례에 따름)
\begin{framedtext}
ㄱ. 乙의 甲에 대한, 매매를 원인으로 하는 소유권이전등기청구권은 소멸시효에 걸리지 않는다. \par
ㄴ. 丙이 토지를 점유ㆍ사용하는 동안에는 丙의 乙에 대한, 매매를 원인으로 하는 소유권이전등기청구권은 소멸시효에 걸리지 않는다. \par
ㄷ.만약 丁이 乙의 점유를 침탈했더라도, 乙의 甲에 대한, 매매를 원인으로 하는 소유권이전등기청구권은 소멸시효가 진행하지 않는다. \par
ㄹ.만약 2014년 4월 戊가 丙의 점유를 침탈했다면, 2015년 6월 현재 丙은 戊에게 점유물반환청구를 할 수 있다.
\end{framedtext}
\begin{choices}
  \item ㄱ, ㄴ
  \item ㄱ, ㄹ
  \item ㄴ, ㄷ
  \item ㄷ, ㄹ
  \item ㄱ, ㄴ, ㄷ
\end{choices}
\end{question}
\answer{1}
